%
%============= CHAPTER BREAK =================%
%
\chapter{Galen (ca 168--192 CE)}%
\vspace*{3mm}%
\justifying%
%
\noindent{}The most prolific of any ancient Greek author, %
Galen left behind some four million words of extant works.\endnote{%
The word count comes from the Stanford Encyclopedia of Philosophy article on Galen: %
\url{plato.stanford.edu/entries/galen/}} %
Born in 129 CE in Pergamon, Asia Minor, he moved to Rome around 159 CE, %
where he did almost all of his writing. %
The personal physician to Emperor Commodus, %
he was famous for his medical acumen and infamous for his contentious mannerisms. %
\vfill{}
\noindent{}\textbf{ἐπιλείχω} / \textit{epileichō} / “to lick over” %
only appears four times in the TLG %
and does not seem to appear in any extant Greek papyri or inscriptions. %
Its first non-Christian occurrence in the TLG is found in the voluminous writings of Galen, %
who uses the term just once. %
In his work \textit{Commentary on Hippocratesʼ on Fractures} / \textit{In Hippocratis librum de fracturis}, %
Galen describes the procedure for applying a bandage in the area of the spleen: %
\vfill{}
 \begin{quote}ὑπὸ δυοῖν ἀρχῶν ἐπίδεσμος ἐπιδεῖται, %
 τουτέστι τῆς μεσότητος τῶν σπληνῶν ... % 
 ἐγκλινομένων ὀλίγον ὡς ἀλλήλων συμβαλεῖν κατὰ τὸ ἀντικείμενον μέρος %
 ἢ ὅθεν ἤρξατο καὶ γενέσθαι τὴν συμβολὴν αὐτῷ τῷ χ γράμματι παραπλησίαν, % 
 ὅπερ εἰ κυκλοτερῶς \textbf{ἐπιλειχθήσεται}, οὐκ ἂν ἐγένετο. %
 συμβαίνει γὰρ τοσούτῳ περιλαμβανομένου ἀπαντᾷν ἀλλήλοις ἄντικρυς. (Kühn 18.2:564) \\[1mm] %
 the bandage is fastened by two starting points, %
 that is from the middle of the spleen with the bandage tails ... %
 inclining a little so that they connect with each other on the opposite side %
 from which it began, and with it making a symbol resembling the letter χ, %
 which if it \textbf{were licked over} in a circular shape, would not work. %
 For in such manner each joins, being wrapped to meet opposite to each other. %
 \end{quote} % 
\clearpage

The apparently new lemma appears in canonical Luke 16:21, %
in the story of the Rich Man and Lazarus. %
The vivid description of dogs licking the beggar was most likely not present in Marcionʼs \textit{Evangelion}, %
particularly since it goes unmentioned by Tertullian and Epiphanius, %
who otherwise confirm many of the basic elements of the story.\endnote{%
It is possibly attested here in the \textit{Adamantius Dialogue}, book 5, %
but scholars have often challenged the reliability of that work and particularly called into %
question the reliability of its attestations to Marcionʼs \textit{Evangelion}, %
especially book 5.} %

\vspace*{2mm} 
 \begin{quote}πτωχὸς δέ τις ὀνόματι Λάζαρος ἐβέβλητο πρὸς τὸν πυλῶνα αὐτοῦ εἱλκωμένος %
 ²¹ καὶ ἐπιθυμῶν χορτασθῆναι ἀπὸ τῶν πιπτόντων ἀπὸ τῆς τραπέζης τοῦ πλουσίου· %
 ἀλλὰ καὶ οἱ κύνες ἐρχόμενοι \textbf{ἐπέλειχον} τὰ ἕλκη αὐτοῦ (Lk 16:20--21) \\[1mm] %
 But a certain beggar, Lazarus by name, %
 had been cast upon his gate, wounded, %
 ²¹ and longing to be filled from what fell from the table of the rich man. %
 But the dogs, coming over, \textbf{were licking over} his wounds. %
 \end{quote} % 
 \vspace*{2mm}

It seems unlikely to be mere coincidence that a new, rare lemma %
appears alongside a vivid depiction of a wound-afflicted person. %
As many \textit{nomina sacra} --- the earliest form of Christian art, %
evident in many early Gospel manuscripts --- would indicate, %
early Christians were fascinated with the letter χ as a symbol, %
both for “Christ” / Χριστός and the cross. %
This passage in Galen, referencing the formation of a χ shape as the means of treating an internal wound, %
would have been a prime candidate for Christian use and comment. %
This points to the canonical redactor of Luke, most likely sometime after the 160s CE, %
drawing on Galenʼs medical commentary to make the portrayal of the beggar Lazarus %
(the Greek form of the Hebrew name Eleazar) that much more vivid. %
\clearpage

\noindent{}\textbf{μυλικός} / \textit{mulikos} / “pestle|grinding” only appears %
one time in the NT, in canonical Luke 17:2, %
but apparently not in its attested parallel in Marcionʼs \textit{Evangelion}, %
canonical Mark, or canonical Matthew, %
which all apparently reflect an earlier formulation. %
Compare:

\vspace*{2mm} 
 \begin{quote}συμφέρει αὐτῷ εἰ μὴ ἐγεννήθη %
    ἢ γεννηθέντα μύλῳ ὀνικῷ προστεθῆναι περὶ τὸν τράχηλον αὐτοῦ 
 καὶ καταποντισθῆναι ἐν τῷ βάθει τῆς θαλάσσης (*Ev 17:2)\endnote{%
 For attestations and critical reconstruction, see %
 Bilby 2020/2025 vol. 2 at Lk2 17.1--2.} \\[1mm] %
It is beneficial for him had he not been born, or having been born, %
 for a donkey millstone to be placed upon his neck, %
 and to be thrown down a cliff into the depth of the sea. %
 \end{quote} % 
 \vspace*{1mm}

 \begin{quote} καλόν ἐστιν αὐτῷ μᾶλλον εἰ περίκειται μύλος ὀνικὸς περὶ τὸν %
 τράχηλον αὐτοῦ καὶ βέβληται εἰς τὴν θάλασσαν. (Mark 9:42) \\[1mm] %
 It is good for him instead if a donkey millstone is put around %
 his neck and he is cast into the sea. \end{quote} % 
 \vspace*{1mm}

 \begin{quote} %
 συμφέρει αὐτῷ ἵνα κρεμασθῇ μύλος ὀνικὸς περὶ τὸν %
τράχηλον αὐτοῦ καὶ καταποντισθῇ ἐν τῷ πελάγει τῆς θαλάσσης. (Matt 18:6)\endnote{%
Note also that the saying in Ev 17:2 that “had he not been born” %
is separated out and attached to Judas in canonical Matthew 26:24. %
} \\[1mm] %
It is beneficial for him if a donkey millstone be hung around his neck %
and he were thrown down a cliff into the high sea. %
\end{quote} %
\vspace*{1mm}

 \begin{quote}λυσιτελεῖ αὐτῷ εἰ λίθος \textbf{μυλικὸς} περίκειται περὶ τὸν τράχηλον αὐτοῦ %
 καὶ ἔρριπται εἰς τὴν θάλασσαν. (Lk 17:2) \\[1mm] %
 It is profitable for him if a \textbf{grinding} stone is wrapped around his neck %
 and he is tossed into the sea. \end{quote} % 
\vspace*{2mm}

\noindent{}The first three passages run more closely to each other %
and to the famous death of Aesop %
as narrated in the \textit{Life of Aesop}, where the %
protagonist is hurled off a cliff to his death. %
Canonical Luke is the outlier in several respects, %
especially the substitution of a “grinding stone” in place of a “donkey millstone”.\endnote{%
See the concluding chapter for ὀνικός / \textit{onikos} / “of, for, or like a donkey” as a %
false positive post-Hadrianic neologism. %
In the Gospels, its use likely derived from the primary source behind Marcionʼs \textit{Evangelion}.} %

The first chronologically secure TLG hit for μυλικός %
belongs to Galen, who uses the lemma seven (!) different times %
in his treatise \textit{De compositione medicamentorum} / %
\textit{On the Composition of Medicines}. %

\vspace*{3mm}
\begin{quote}
[\textbf{Μυλικὴ} ἀγαθὴ, ἐπὶ πολλῶν θαυμασθεῖσα.] (\textit{C.M.} 12:869) \\[0.5mm] %
[A good \textbf{grinding}, admired for many uses.]
\end{quote} %
\vspace*{1mm}

\begin{quote}
[Τὰ ὑπ’ Ἀνδρομάχου γεγραμμένα πρὸς ἁπάσας ὀδύνας ὀδόντων %
\textbf{μυλικὴ} αὐθωρὸν παύουσα.] (\textit{C.M.} 12:877) \\[0.5mm]
[The things prescribed by Andromachus for all toothaches, %
a \textbf{grinding} that immediately stops [pain].]
\end{quote} %
\vspace*{1mm}

\begin{quote}
ἄλλη \textbf{μυλικὴ} ἐκ τῶν Ἀφρόδα. ... ἄλλη \textbf{μυλικὴ} τοῦ αὐτοῦ· ... %
\textbf{μυλικὴ} καλή. (\textit{C.M.} 12:878) \\[0.5mm]
Another \textbf{grinding} from Aphrodas. ... Another \textbf{grinding} from him. ... A good \textbf{grinding}. % 
\end{quote} %
\vspace*{1mm}

\begin{quote}
\textbf{μυλικὴ} ἄλλη ὡς Ἀριστοκράτης. (\textit{C.M.} 12:879) \\[0.5mm]
Another \textbf{grinding}, as by Aristocrates.
\end{quote} %
\vspace*{1mm}

\begin{quote}
Ἄλλη \textbf{μυλικὴ} αὐθωρὸν παύουσα. (\textit{C.M.} 12:882) \\[0.5mm]
Another \textbf{grinding}, quickly stopping [pain].
\end{quote}
\vspace*{3mm}

\noindent{}If we exclude the bracketed notes, which may be supplemental annotations, %
five hits remain. Galen also uses the term μυλητικὴ / %
“a ground [paste]” (\textit{C.M.} 12:877) %
as a topical dental anesthetic.

This lemma does not seem to appear among digitized Greek papyri or Greek inscriptions. %
The Lukan substitution is likely indebted to Galen in some form or fashion. %
It essentially re-purposes a newly popularized medicinal lemma associated %
with the concoction of ground pastes. %
In doing so, it apparently pictures a smaller, more readily available type of stone %
than that described in Marcionʼs \textit{Evangelion}, canonical Mark, or canonical Matthew. % 
One also wonders if the substitution befits the sea-faring obsession of the editor(s) %
of Luke-Acts: painting a plausible scenario of execution at sea with a smaller weight %
strapped to the victimʼs neck (perhaps implying a shipʼs doctor), %
rather than a land-based setting that would require a beast of burden to transport %
an inhuman weight of execution. %
%